%% kapa: exam=1 round=36 year=2025 period=1 form=A subject=civil_law subject_ko=민법 seq=39 number=39
%% source: Q-Net 감정평가사 기출문제 (cst00309) artlSeq=5231525 file=raw/1008/5231525/2025년 제36회 감정평가사 1차시험 1교시 문제지 원본.pdf
%% conversion: kordoc 4.24.51 -> tools/gen1.py (auto); review=proofread-sonnet; flags=none
\begin{question}{39}
저당권에 관한 설명으로 옳은 것을 모두 고른 것은? (다툼이 있으면 판례에 따름)
\begin{framedtext}
ㄱ. 건물에 저당권이 설정된 후 건물의 종물에 대해 강제집행을 한 자는 건물 경매절차의 매수인에게 강제집행의 효력을 주장할 수 없다. \par
ㄴ. 저당권이 설정된 토지에 지상권을 취득한 자가 수목을 식재한 경우 그 지상권자는 토지 경매절차의 매수인에 대해 수목의 매수를 청구할 수 있다. \par
ㄷ. 건물에 설정된 저당권이 실행된 경우 경매절차의 매수인은 특별한 사정이 없는 한 건물의 소유를 위해 인정된 지상권도 등기 없이 취득한다.
\end{framedtext}
\begin{choices}
  \item ㄱ
  \item ㄴ
  \item ㄱ, ㄷ
  \item ㄴ, ㄷ
  \item ㄱ, ㄴ, ㄷ
\end{choices}
\end{question}
\answer{3}
